\documentclass[UTF8,a4paper,11pt]{ctexart}

\usepackage{geometry}
\usepackage{booktabs}
\usepackage{longtable}
\usepackage{array}
\usepackage{ragged2e}
\usepackage{xcolor}
\usepackage{graphicx}
\usepackage{float}
\usepackage{tikz}
\usetikzlibrary{arrows.meta,positioning,calc}
\PassOptionsToPackage{obeyspaces}{url}
\usepackage{xurl}
\usepackage{hyperref}
\usepackage{fancyhdr}
\usepackage{lastpage}

\geometry{left=2.0cm,right=2.0cm,top=2.2cm,bottom=2.2cm}
\hypersetup{colorlinks=true,linkcolor=blue,urlcolor=blue,citecolor=blue}
\pagestyle{fancy}
\fancyhf{}
\lhead{委员会报告}
\rhead{\thepage/\pageref{LastPage}}
\setlength{\headheight}{14pt}
\setlength{\parindent}{2em}
\setlength{\parskip}{0.42em}
\setlength{\tabcolsep}{2.2pt}
\renewcommand{\arraystretch}{1.15}
\emergencystretch=3em

\newcolumntype{L}[1]{>{\RaggedRight\arraybackslash}p{#1}}
\newcolumntype{C}[1]{>{\Centering\arraybackslash}p{#1}}
\DeclareUrlCommand\code{\urlstyle{tt}}
\newcommand{\wf}[1]{\texttt{\detokenize{#1}}}
\newcommand{\inst}[1]{\texttt{\detokenize{#1}}}
\newcommand{\pass}{\textcolor{green!45!black}{通过}}
\newcommand{\fail}{\textcolor{red!65!black}{未通过}}
\newcommand{\uneval}{\textcolor{orange!75!black}{未评测}}

\tikzset{
  box/.style={
    rectangle,
    rounded corners=2pt,
    draw=black!55,
    line width=0.35pt,
    align=center,
    minimum height=8mm,
    inner xsep=3pt,
    inner ysep=3pt,
    font=\small
  },
  evidence/.style={box, fill=cyan!9, draw=cyan!45!black},
  judge/.style={box, fill=orange!12, draw=orange!60!black},
  coder/.style={box, fill=green!10, draw=green!45!black},
  risk/.style={box, fill=red!8, draw=red!55!black},
  final/.style={box, fill=violet!9, draw=violet!45!black},
  arrow/.style={-{Latex[length=2.2mm,width=1.5mm]}, line width=0.45pt, draw=black!70},
  softarrow/.style={-{Latex[length=2.0mm,width=1.3mm]}, line width=0.4pt, draw=black!45, dashed}
}

\title{\textbf{委员会报告}\\OpenCollab + GLM-5.2 的 SWE 委员会原理与实验}
\author{实验整理}
\date{2026-06-30}

\begin{document}
\maketitle

\begin{abstract}
这份报告把 OpenCollab 中“委员会模式”的设计动机、工作流结构和 SWE-bench Lite 前 100 题上的补测结果单独整理出来。委员会模式的核心思想是把软件修复任务拆成定位、证据、验证、修复、补丁攻击和最终放行几个岗位，让模型在写生产补丁前先建立可审计证据，在写完补丁后再接受边界、回归和可见行为变化审查。报告特别展开这些岗位的具体输入、动作、产物和判断标准，避免把“契约”“验证”“攻击”停留在抽象名词上。

当前已经进入实验统计的 SWE 委员会补测使用 V1 工作流，代码名为 \code{validation-council-solve}。

\end{abstract}

\section{问题背景}

SWE-bench 的任务表面上是“给定一个 issue，让模型改代码”，实际难点在于隐藏测试检查的是更精确的行为契约。这里的“行为契约”指可被程序观察到的要求，例如某个函数应返回什么值、应抛出哪类异常、是否应发出 warning、输出顺序是否稳定、迁移前后数据库里应该留下什么记录、某个选项是否应该被忽略。模型在生成补丁时只能看到问题文本、起始版本仓库、公开测试和公开文档。官方隐藏测试、官方评分补丁、FAIL\_TO\_PASS 测试节点在生成阶段都不可见，只能在最终评分阶段使用。

这里有一个直接的矛盾。修复者如果过早写代码，很容易把某个合理猜测写成生产行为。测试者如果缺少证据分层，也可能把没有来源的猜测写成临时测试预期。委员会模式尝试把这两个风险拆开处理。它要求不同角色先回答“证据支持什么”，再回答“代码应该怎么改”，最后回答“补丁有没有改出新的可见行为风险”。

\begin{figure}[H]
\centering
\resizebox{0.94\textwidth}{!}{
\begin{tikzpicture}[node distance=9mm and 12mm]
\node[evidence, text width=2.0cm] (issue) {问题文本\\缺陷描述};
\node[evidence, text width=2.0cm, below=of issue] (repo) {起始版本\\源码};
\node[evidence, text width=2.0cm, below=of repo] (public) {公开测试\\公开文档};
\node[coder, text width=3.0cm, right=18mm of repo] (workflow) {OpenCollab\\生成阶段};
\node[final, text width=1.8cm, right=of workflow] (patch) {模型\\补丁};
\node[judge, text width=3.0cm, right=of patch] (harness) {官方评测器\\评分阶段};
\node[final, text width=1.8cm, right=of harness] (score) {是否\\通过};
\node[risk, text width=2.6cm, below=13mm of harness] (hidden) {隐藏测试\\评分答案};
\draw[arrow] (issue.east) -- (workflow.west);
\draw[arrow] (repo.east) -- (workflow.west);
\draw[arrow] (public.east) -- (workflow.west);
\draw[arrow] (workflow) -- (patch);
\draw[arrow] (patch) -- (harness);
\draw[arrow] (harness) -- (score);
\draw[arrow] (hidden) -- node[right, font=\scriptsize] {最终评分使用} (harness);
\draw[softarrow] (hidden.west) to[out=180,in=-55] node[below, font=\scriptsize] {生成阶段不可见} (workflow.south east);
\end{tikzpicture}
}
\caption{SWE-bench 盲验证口径下，模型生成阶段和官方评分阶段的边界}
\end{figure}

图里最关键的一条边界，是委员会可以构造自己的临时验证，但这些验证必须来自公开信息。所谓“验证更强”，指的是模型和工作流从公开信息里挖出更多行为约束，而不是提前看到官方隐藏测试。

\section{从团队到委员会}

OpenCollab 原有团队模式依赖负责人根据现场情况调度角色。负责人可以安排测试员、修复者和评审者，也可以根据中间结果决定继续探索或停止。这样的方式灵活，但重复实验时控制流会随负责人判断变化，某些关键步骤也可能被跳过。

工作流模式把控制权放进 Python 代码。节点顺序、并行扇出、最大重试次数、放行条件和失败分支都由代码固定，每个节点内部仍由模型完成推理和文本生成。委员会模式就是这种思路在 SWE-bench 上的专门版本。它保留多角色协作，但把岗位关系和状态流转固定成可审计的实验协议。

\begin{figure}[H]
\centering
\resizebox{\textwidth}{!}{
\begin{tikzpicture}[node distance=8mm and 9mm]
\node[final, text width=1.8cm] (teamissue) {问题};
\node[coder, text width=2.0cm, right=of teamissue] (lead) {负责人\\临场调度};
\node[evidence, text width=1.7cm, above right=5mm and 8mm of lead] (tester) {测试员};
\node[coder, text width=1.7cm, right=of lead] (coder1) {修复者};
\node[judge, text width=1.7cm, below right=5mm and 8mm of lead] (reviewer) {评审者};
\node[final, text width=1.8cm, right=19mm of coder1] (teamout) {补丁};
\draw[arrow] (teamissue) -- (lead);
\draw[arrow] (lead) -- (tester);
\draw[arrow] (lead) -- (coder1);
\draw[arrow] (lead) -- (reviewer);
\draw[arrow] (tester) -- (teamout);
\draw[arrow] (coder1) -- (teamout);
\draw[arrow] (reviewer) -- (teamout);
\node[final, text width=1.8cm, right=23mm of teamout] (wfissue) {问题};
\node[evidence, text width=2.0cm, right=of wfissue] (evi) {证据节点};
\node[judge, text width=2.0cm, right=of evi] (j1) {验证裁判};
\node[coder, text width=1.7cm, right=of j1] (wcoder) {修复者};
\node[risk, text width=2.0cm, right=of wcoder] (attack) {补丁攻击};
\node[final, text width=1.8cm, right=of attack] (wfout) {补丁};
\draw[arrow] (wfissue) -- (evi);
\draw[arrow] (evi) -- (j1);
\draw[arrow] (j1) -- (wcoder);
\draw[arrow] (wcoder) -- (attack);
\draw[arrow] (attack) -- (wfout);
\draw[softarrow] (attack.south) to[out=-90,in=-90,looseness=0.8] node[below, font=\scriptsize] {失败分支由代码固定} (j1.south);
\node[font=\small, above=8mm of lead] {团队模式};
\node[font=\small, above=8mm of j1] {委员会工作流};
\end{tikzpicture}
}
\caption{团队模式与委员会工作流的控制权差异}
\end{figure}

这种结构带来的科研价值在于可复查。若某个补丁失败，报告可以追问失败来自契约证据不足、边界验证不足、回归审查不足，还是修复者实现错误。若某个补丁成功，也可以观察成功来自更准的定位、更强的验证，还是补丁攻击阶段拦下了错误方向。

\section{术语拆解：这些岗位具体做什么}

“契约提取”听起来抽象，实际动作很具体。契约挖掘者会先读 issue 原文、issue 讨论中公开的复现描述、相关源码、公开测试和公开文档，然后把它们改写成若干条可观察行为。每条行为都要说明触发条件、观察对象、期望结果和证据来源。例如“当自引用对称 ManyToManyField 显式传入 \code{related_name} 时，应出现 \code{fields.W345} warning”就是一条契约。反过来，“所有 ManyToManyField 都要出现这个 warning”就缺少证据，应该被降级为猜测或直接拒绝。

“测试地图”负责把一个行为要求翻译成仓库内可执行的验证位置。测试地图员要回答的是这个仓库里到底怎样验证一个行为，包括测试命令是什么、已有测试放在哪些目录、测试夹具如何初始化、公开测试里有没有相近断言、临时 probe 应该放在哪里、运行后怎样清理、哪些文件不能进入最终 diff。对 Django 任务来说，它通常会定位到 \code{tests/runtests.py}、某个 \code{tests/...} 目录、model check 的断言风格，以及可用的最小 Python probe。

“验证”分成生成和裁判两步。验证生成者把契约改成候选检查，例如一段临时 Python 脚本、一个公开测试命令、一个 grep 文档条目的命令。验证裁判只保留证据足够的检查。证据强的检查可以写成硬预期，例如“这三种输入分别应产生 warning、无 warning、无 warning”。证据弱的检查只能做诊断，例如“这个边界值得观察，但不能作为失败标准”。这样可以减少模型把猜测写成测试答案的风险。

“补丁攻击”发生在修复者已经改完代码之后，审查对象是当前 diff 背后的行为假设。分支攻击会改变输入组合，看补丁是否只覆盖了最显眼的分支。回归扫描会找相邻公开行为，看补丁有没有破坏旧功能。可见行为审查会检查返回值、异常、warning 文本、默认值、顺序和文档条目是否被证据支持。攻击阶段最后产出一组需要修复者处理的具体风险。

\begin{center}
{\small
\setlength{\tabcolsep}{3pt}
\begin{longtable}{L{3.0cm} L{5.4cm} L{5.8cm}}
\toprule
术语 & 具体动作 & 合格产物长什么样 \\
\midrule
\endfirsthead
\toprule
术语 & 具体动作 & 合格产物长什么样 \\
\midrule
\endhead
定位 & 找到最可能相关的文件、函数、类和调用路径，写出根因假设。 & 指向可读源码位置，并说明为什么这些位置和 issue 症状有关。 \\
契约提取 & 把公开材料改写成可观察行为，标注触发条件、观察对象、期望结果和证据来源。 & 每条契约都能追到 issue、源码、公开测试或文档，缺证据的内容被标为猜测。 \\
测试地图 & 找测试框架、测试命令、夹具、相邻测试、临时 probe 写法和清理规则。 & 给出可运行命令和最小验证方式，同时说明临时文件不进入最终 diff。 \\
修复前验证 & 在起始版本上运行或规划检查，区分 base-fail、base-pass、invalid、weak。 & 修复者知道哪些检查复现 bug，哪些只是诊断，哪些检查无效。 \\
修复 & 在生产源码中做最小改动，避免改测试或提交临时文件。 & diff 只包含与 issue 根因相关的源码或文档改动。 \\
补丁攻击 & 围绕当前 diff 设计边界输入、相邻行为检查和可见行为审查。 & 输出具体风险，例如“这个 warning 是否只应在自引用字段触发”。 \\
最终验证 & 读取最终 diff、运行可行测试、确认临时文件缺席、列出允许提交的路径。 & 给出 PASS/FAIL、允许路径、拒绝路径和失败理由。 \\
\bottomrule
\caption{委员会术语的可执行含义}
\end{longtable}
}
\end{center}

\section{委员会 V1：实际运行的验证委员会}

当前已经跑过并进入统计的委员会实验使用 V1，代码名是 \wf{validation-council-solve}。所有 SWE 委员会补测脚本传入的工作流参数都是 \wf{--workflow validation-council-solve}，并开启盲验证。V1 的设计目标是把“写代码”和“验证代码”拆开。修复者写补丁前要先接收契约和测试地图，写完补丁后要经过补丁验证、风险审计、修复后验证和最终验证。

\begin{figure}[H]
\centering
\resizebox{\textwidth}{!}{
\begin{tikzpicture}[node distance=7mm and 8mm]
\node[final, text width=1.5cm] (issue) {问题};
\node[evidence, text width=1.8cm, right=of issue] (loc) {定位者};
\node[evidence, text width=2.0cm, above right=5mm and 7mm of loc] (cm) {契约\\挖掘者};
\node[evidence, text width=2.0cm, below right=5mm and 7mm of loc] (tc) {测试\\地图员};
\node[judge, text width=2.0cm, right=21mm of loc] (pre) {修复前\\验证};
\node[judge, text width=1.8cm, right=of pre] (base) {起始版\\分诊};
\node[coder, text width=1.6cm, right=of base] (coder) {修复者};
\node[judge, text width=1.9cm, right=of coder] (pv) {补丁\\验证};
\node[risk, text width=2.0cm, right=of pv] (risk) {补丁风险\\审计};
\node[judge, text width=2.0cm, right=of risk] (post) {修复后\\验证};
\node[final, text width=1.8cm, right=of post] (fv) {最终\\验证};
\node[final, text width=1.5cm, right=of fv] (out) {补丁};
\draw[arrow] (issue) -- (loc);
\draw[arrow] (loc) -- (cm);
\draw[arrow] (loc) -- (tc);
\draw[arrow] (cm) -- (pre);
\draw[arrow] (tc) -- (pre);
\draw[arrow] (pre) -- (base);
\draw[arrow] (base) -- (coder);
\draw[arrow] (coder) -- (pv);
\draw[arrow] (pv) -- (risk);
\draw[arrow] (risk) -- (post);
\draw[arrow] (post) -- (fv);
\draw[arrow] (fv) -- (out);
\draw[softarrow] (fv.south) to[out=-95,in=-90,looseness=0.68] node[below, font=\scriptsize] {失败后带反馈小修} (coder.south);
\end{tikzpicture}
}
\caption{委员会 V1：验证委员会的实际运行结构}
\end{figure}

V1 的前半段负责建立证据。定位者先把 issue 的自然语言症状变成源码范围，例如“相关逻辑可能在某个字段检查函数、迁移函数或渲染函数附近”。契约挖掘者接着把问题拆成期望行为、当前错误行为和不应改变的旧行为。测试地图员补充仓库层面的验证信息，说明应该用哪个测试命令、哪些公开测试最接近、临时 probe 怎样写才不会污染最终 diff。随后，修复前验证工厂把这些材料改成候选检查，修复前验证裁判最多保留 5 个有证据支撑的候选，起始版分诊器再判断这些候选在原始版本上的状态。

V1 的后半段负责审查补丁。修复者最多进行 2 轮尝试，每轮之后由补丁验证者检查差异边界，补丁风险审计者寻找边界和回归风险，修复后验证工厂和裁判生成并筛选新的检查，最终验证者读取差异、运行可行的公开测试和已批准验证，并给出允许进入模型补丁的路径。最终验证者通过时，工作流产出 guarded staged diff。最终验证者失败时，失败摘要会反馈给下一轮修复者。

换成一次真实运行来看，V1 的每一步都对应一个可检查的中间产物。以 \inst{django__django-14730} 为例，这个 issue 说的是自引用对称 ManyToManyField 上 \code{related_name} 实际无效，但 Django 当时没有给出任何提示。定位者把范围缩到 \code{django/db/models/fields/related.py} 的 ManyToManyField 检查逻辑，以及 \code{docs/ref/checks.txt} 的检查项文档。契约挖掘者从中抽出三条核心行为。显式传入 \code{related_name} 且关系是自引用对称时应给 warning，没有传入 \code{related_name} 时不应给这个 warning，已有的 W340、W343 等检查则要保持原样。测试地图员接着指出可用的验证方式，包括构造最小 Django model、调用字段检查函数、运行相关 model check 公开测试，再检查文档里是否补充了 W345。

后面的角色继续把这些产物接上。修复前验证裁判会接受能区分三种输入的临时 probe，也会接受公开测试命令和文档 grep。修复者据此保存用户原始传入的 \code{related_name}，在检查函数里新增 \code{fields.W345} warning，并补上文档条目。补丁风险审计者再看这个 diff 是否把 warning 范围放得过宽，是否误伤非自引用字段，是否把 warning 变成 error，是否漏了文档。最终验证者只允许 \code{django/db/models/fields/related.py} 和 \code{docs/ref/checks.txt} 进入预测补丁，并确认临时 probe 没有留在 diff 里。这个例子说明，“委员会”真正做的是把一条 issue 逐步转成一组可观察行为、一组验证命令和一个受约束的生产 diff。

\begin{center}
{\small
\setlength{\tabcolsep}{3pt}
\begin{tabular}{L{2.8cm} L{4.0cm} L{4.0cm} L{3.3cm}}
\toprule
角色 & 主要输入 & 具体动作 & 主要产物 \\
\midrule
定位者 & 问题文本、源码、公开测试 & 找相关文件和根因路径，写完成条件。 & 根因假设、相关文件、完成条件 \\
契约挖掘者 & 定位结果、公开信息 & 把自然语言症状改写成可观察行为。 & 行为契约、无影响行为、证据来源 \\
测试地图员 & 仓库测试结构 & 找测试命令、夹具、相邻测试和临时 probe 写法。 & 可运行命令、测试入口、临时验证建议 \\
修复前验证裁判 & 候选验证、行为契约 & 接受强证据检查，拒绝猜测预期，保留诊断信号。 & 被批准的临时检查、拒绝理由 \\
修复者 & 契约、测试地图、分诊结果、失败反馈 & 只改生产源码或必要文档，保持 diff 最小。 & 最小生产补丁 \\
补丁风险审计者 & 当前差异、契约、测试地图 & 找补丁覆盖不到的分支、边界和相邻回归。 & 边界风险、回归风险、需补充验证的字段 \\
最终验证者 & 当前差异、所有验证结果 & 查 diff、跑可行测试、确认临时文件缺席。 & 放行判定、允许提交路径 \\
\bottomrule
\end{tabular}
}
\end{center}

\section{委员会 V2：后续图同构设计}

V2 的代码名是 \wf{swe-committee-v2}。它来自 V1 的经验，但结构更细，目标是让证据阶段、仲裁阶段、补丁攻击阶段和最终放行阶段严格分工。V2 增加了可见行为清单员和契约仲裁庭。可见行为清单员会把返回值、异常类型、warning 文本、错误码、默认值、排序、数据库副作用、文档承诺等外部可观察字段列出来。契约仲裁庭再把证据分成强证据、弱证据、猜测字段和禁止硬断言字段。强证据可以生成硬测试，弱证据只能生成诊断，猜测字段只保留为假设，禁止字段会被后续最终怀疑者重点检查。

\begin{figure}[H]
\centering
\resizebox{0.92\textwidth}{!}{
\begin{tikzpicture}[node distance=8mm and 12mm]
\node[evidence, text width=2.1cm] (cm) {契约\\挖掘者};
\node[evidence, text width=2.1cm, below=of cm] (tc) {测试\\地图员};
\node[evidence, text width=2.3cm, below=of tc] (ci) {可见行为\\清单员};
\node[judge, text width=2.2cm, right=18mm of tc] (tribunal) {契约\\仲裁庭};
\node[coder, text width=2.4cm, above right=1mm and 17mm of tribunal] (strong) {强证据\\硬测试};
\node[judge, text width=2.4cm, right=17mm of tribunal] (weak) {弱证据\\诊断};
\node[final, text width=2.4cm, below right=1mm and 17mm of tribunal] (guess) {猜测字段\\待确认};
\node[risk, text width=2.4cm, below=of guess] (forbid) {禁止字段\\不可硬断言};
\draw[arrow] (cm) -- (tribunal);
\draw[arrow] (tc) -- (tribunal);
\draw[arrow] (ci) -- (tribunal);
\draw[arrow] (tribunal) -- (strong);
\draw[arrow] (tribunal) -- (weak);
\draw[arrow] (tribunal) -- (guess);
\draw[arrow] (tribunal) -- (forbid);
\end{tikzpicture}
}
\caption{委员会 V2 的证据分层}
\end{figure}

V2 的补丁后阶段也比 V1 更明确。已有测试和批准验证先检查补丁状态，随后分支攻击者寻找条件分支和边界输入，回归扫描者关注相邻行为是否被破坏，可见变化审查者检查返回值、异常、警告、文本、默认值和顺序等外部可见字段是否被证据授权。修复后验证裁判看到有证据失败时，直接把任务交给修复者小修。最终怀疑者看到契约证据不足时，任务回到契约仲裁庭。最终验证者看到具体提交问题时，任务回到修复者。

V2 里“强证据”和“弱证据”的差别也可以具体说明。强证据通常来自 issue 明确描述、公开测试已有断言、源码注释或文档承诺，可以写成必须满足的检查。弱证据通常来自相邻实现、合理推断或风险审计，它可以提醒模型观察某个方向，但不能直接写成 expected output。猜测字段表示当前材料无法支持，最终怀疑者会阻止补丁依赖这些字段。比如某个 issue 只说“应该发出 warning”，但没有说明 warning 文本必须逐字相同，那么“有 warning”可能是强证据，“具体英文句子逐字匹配”就需要额外文档或源码证据。

\begin{figure}[H]
\centering
\resizebox{\textwidth}{!}{
\begin{tikzpicture}[node distance=7mm and 8mm]
\node[final, text width=1.4cm] (issue) {问题};
\node[evidence, text width=1.7cm, right=of issue] (a) {定位者};
\node[evidence, text width=1.9cm, above right=3mm and 7mm of a] (cm) {契约\\挖掘};
\node[evidence, text width=1.9cm, right=of a] (tc) {测试\\地图};
\node[evidence, text width=2.1cm, below right=3mm and 7mm of a] (ci) {可见行为\\清单};
\node[judge, text width=2.1cm, right=16mm of tc] (ca) {契约\\仲裁};
\node[judge, text width=2.0cm, right=of ca] (ptd) {修复前\\验证};
\node[judge, text width=1.6cm, right=of ptd] (vj1) {前置\\裁判};
\node[judge, text width=1.7cm, right=of vj1] (be) {起始版\\分诊};
\node[coder, text width=1.5cm, right=of be] (coder) {修复者};
\node[final, text width=2.1cm, right=of coder] (etv) {已有测试\\批准验证};
\node[risk, text width=1.7cm, above right=3mm and 7mm of etv] (da) {分支\\攻击};
\node[risk, text width=1.7cm, right=of etv] (rs) {回归\\扫描};
\node[risk, text width=1.8cm, below right=3mm and 7mm of etv] (odr) {可见变化\\审查};
\node[judge, text width=2.0cm, right=16mm of rs] (post) {修复后\\验证};
\node[judge, text width=1.7cm, right=of post] (vj2) {后置\\裁判};
\draw[arrow] (issue) -- (a);
\draw[arrow] (a) -- (cm);
\draw[arrow] (a) -- (tc);
\draw[arrow] (a) -- (ci);
\draw[arrow] (cm) -- (ca);
\draw[arrow] (tc) -- (ca);
\draw[arrow] (ci) -- (ca);
\draw[arrow] (ca) -- (ptd);
\draw[arrow] (ptd) -- (vj1);
\draw[arrow] (vj1) -- (be);
\draw[arrow] (be) -- (coder);
\draw[arrow] (coder) -- (etv);
\draw[arrow] (etv) -- (da);
\draw[arrow] (etv) -- (rs);
\draw[arrow] (etv) -- (odr);
\draw[arrow] (da) -- (post);
\draw[arrow] (rs) -- (post);
\draw[arrow] (odr) -- (post);
\draw[arrow] (post) -- (vj2);
\end{tikzpicture}
}
\caption{委员会 V2 的主结构：三路证据、前置验证、三路补丁攻击}
\end{figure}

\begin{figure}[H]
\centering
\resizebox{0.88\textwidth}{!}{
\begin{tikzpicture}[node distance=10mm and 14mm]
\node[judge, text width=1.8cm] (vj2) {后置\\裁判};
\node[final, text width=1.8cm, right=of vj2] (fs) {最终\\怀疑者};
\node[final, text width=1.8cm, right=of fs] (fv) {最终\\验证者};
\node[final, text width=1.7cm, right=of fv] (out) {模型\\补丁};
\node[coder, text width=2.0cm, below=16mm of fs] (cr) {修复者\\小修};
\node[final, text width=2.2cm, below=16mm of vj2] (etv) {已有测试\\批准验证};
\node[judge, text width=2.0cm, above=16mm of fs] (ca) {契约\\仲裁庭};
\node[judge, text width=2.6cm, above=16mm of fv] (pre) {修复前验证\\起始版分诊};
\draw[arrow] (vj2) -- node[above, font=\scriptsize] {通过} (fs);
\draw[arrow] (fs) -- node[above, font=\scriptsize] {通过} (fv);
\draw[arrow] (fv) -- node[above, font=\scriptsize] {通过} (out);
\draw[arrow] (vj2.south) to[out=-90,in=180] node[pos=0.35, left, font=\scriptsize] {有证据失败} (cr.west);
\draw[arrow] (fv.south) to[out=-90,in=0] node[pos=0.35, right, font=\scriptsize] {具体阻断} (cr.east);
\draw[arrow] (cr) -- (etv);
\draw[arrow] (etv) -- (vj2);
\draw[arrow] (fs) -- node[left, font=\scriptsize] {证据不足} (ca);
\draw[arrow] (ca) -- (pre);
\draw[arrow] (pre.south) to[out=-90,in=90] (cr.north);
\end{tikzpicture}
}
\caption{委员会 V2 的后半段状态机}
\end{figure}

V1 和 V2 的关系并不复杂。V1 把验证从修复者手中独立出来，V2 再把“验证失败”细分为有证据失败、证据不足和具体提交阻断。V1 更轻，适合快速看清验证工作流本身能带来多少收益。V2 更重，适合研究更强的 validation data 能不能减少那些只差一两个 corner case 的失败。

\begin{center}
{\small
\setlength{\tabcolsep}{3pt}
\begin{tabular}{L{3.0cm} L{5.3cm} L{6.0cm}}
\toprule
维度 & 委员会 V1 & 委员会 V2 \\
\midrule
前置证据 & 契约挖掘者和测试地图员并行产出证据。 & 契约挖掘者、测试地图员、可见行为清单员三路并行。 \\
证据处理 & 契约集合直接供验证生成者和修复者使用。 & 契约仲裁庭将证据分成强证据、弱证据、猜测字段、禁止字段。 \\
补丁后审查 & 补丁验证、风险审计、修复后验证、最终验证线性串联。 & 分支攻击、回归扫描、可见变化审查并行进入后置裁判。 \\
失败分流 & 最终验证失败后带反馈重试，最多 2 轮修复。 & 有证据失败回小修，证据不足回仲裁，具体阻断回小修。 \\
实验状态 & 已用于剩余 13 题补测，并进入 89/100 统计。 & 已实现为可调用工作流，后续实验需单独统计。 \\
\bottomrule
\end{tabular}
}
\end{center}

\begin{center}
{\small
\setlength{\tabcolsep}{3pt}
\begin{tabular}{L{3.0cm} L{5.2cm} L{6.1cm}}
\toprule
V2 状态 & 触发条件 & 工作流下一步 \\
\midrule
有证据失败 & 已批准检查在补丁后仍失败，且检查的预期来自强证据。 & 交给修复者小修，反馈具体失败输入和观察结果。 \\
证据不足 & 补丁改变了可见行为，但仲裁庭没有给出足够证据。 & 回到契约仲裁庭，重新判定该字段能否写成硬预期。 \\
具体阻断 & diff 中有临时文件、允许路径不干净、公开测试未通过。 & 交给修复者清理或最小修补。 \\
最终通过 & 强证据检查通过，诊断信号无阻断，diff 只含允许路径。 & 产出模型补丁，交给官方评测器。 \\
\bottomrule
\end{tabular}
}
\end{center}

\section{实验口径}

委员会补测发生在前 100 题已经达到 87/100 之后。此时剩余 13 个样本仍未解决，因而成为委员会 V1 的补测对象。生成阶段使用 GLM-5.2，并通过 \code{swebench/gen_prediction_workflow.py} 调用 \wf{validation-council-solve}。该驱动对 \wf{validation-council-solve} 和 \wf{swe-committee-v2} 默认启用盲验证保护，只向工作流提供公开信息，并用 guarded staged diff 作为最终预测补丁。

主批次预算从 1M token 开始，后续根据预算信号升到 2M 和 4M，针对性补跑则用更高预算处理仍需复核的样本。官方结果只采用 SWE-bench harness 生成的评测报告。如果某个样本没有形成干净的非空 patch，实验表里记为未评测，保守累计分数把未评测样本一律算作尚未解决。

\section{总体结果}

委员会 V1 在剩余 13 个样本上形成了 10 个可评测非空 patch，其中 2 个通过、8 个未通过，另外 3 个样本没有干净非空 patch 进入正式统计。新增通过样本是 \inst{django__django-11019} 和 \inst{django__django-14730}。这样一来，前 100 题保守累计结果从 87/100 提升到 89/100。如果只看已有官方报告的 97 个样本，通过数为 89/97。

\begin{table}[H]
\centering
{\small
\begin{tabular}{L{5.6cm} C{2.4cm} C{2.4cm} C{2.4cm}}
\toprule
口径 & 已通过 & 未通过 & 未评测 \\
\midrule
委员会补测的 13 个剩余样本 & 2 & 8 & 3 \\
前 100 题累计，未评测按未解决计 & 89 & 11 & 0 \\
前 100 题累计，仅看已有官方报告样本 & 89 & 8 & 3 \\
\bottomrule
\end{tabular}
}
\caption{委员会 V1 补测后的总体结果}
\end{table}

\begin{center}
{\scriptsize
\begin{longtable}{L{4.0cm} C{1.2cm} C{1.0cm} C{1.0cm} C{1.0cm} C{1.2cm} L{4.0cm}}
\toprule
实例 & 状态 & apply & F2P失败 & P2P失败 & patch字符 & 首个失败信号或备注 \\
\midrule
\endfirsthead
\toprule
实例 & 状态 & apply & F2P失败 & P2P失败 & patch字符 & 首个失败信号或备注 \\
\midrule
\endhead
\inst{astropy__astropy-14365} & \fail & 是 & 1 & 0 & 475 & \code{test_roundtrip[True]} \\
\inst{astropy__astropy-7746} & \fail & 是 & 1 & 0 & 995 & \code{test_zero_size_input} \\
\inst{django__django-11019} & \pass & 是 & 0 & 0 & 5063 & 新增通过样本 \\
\inst{django__django-11283} & \uneval & -- & -- & -- & 0 & 没有干净非空 patch \\
\inst{django__django-11564} & \fail & 是 & 2 & 0 & 948 & \code{SCRIPT_NAME} 下 static/media 前缀行为仍不完整 \\
\inst{django__django-11630} & \fail & 是 & 2 & 0 & 1671 & duplicate db table checks 仍有残留失败 \\
\inst{django__django-11905} & \fail & 是 & 2 & 0 & 1119 & \code{isnull} 非布尔值与 iterator 行为仍未覆盖 \\
\inst{django__django-13321} & \fail & 是 & 18 & 0 & 787 & session decode 相关失败较重 \\
\inst{django__django-13768} & \fail & 是 & 1 & 0 & 911 & \code{send_robust} logging 行为仍未满足 \\
\inst{django__django-14155} & \fail & 是 & 3 & 0 & 1792 & \code{ResolverMatch} repr 对 partial 的处理仍不足 \\
\inst{django__django-14730} & \pass & 是 & 0 & 0 & 1699 & 新增通过样本 \\
\inst{django__django-15252} & \uneval & -- & -- & -- & 0 & 干净委员会补测未形成正式非空 patch \\
\inst{django__django-15695} & \uneval & -- & -- & -- & 0 & 没有干净非空 patch \\
\bottomrule
\caption{委员会 V1 补测逐题结果}
\end{longtable}
}
\end{center}

\section{新增通过样本}

\inst{django__django-11019} 是 Django forms media 合并问题。委员会补测最终生成的 patch 修改 \code{django/forms/widgets.py}，把两个列表合并逻辑扩展为多列表合并，并使用稳定拓扑排序保留依赖顺序。在这道题里，“契约”可以拆成两条具体要求。一是多个 media 列表合并后依赖顺序要尽量保留，二是当两个文件的相对顺序冲突时系统仍要发出冲突 warning。测试地图则会关注 Django 里 media 渲染和 widget 相关公开测试怎么跑。这个样本说明，委员会的价值在于修复前契约和补丁后验证一起约束了排序语义，而单纯做局部拼接，很容易遗漏冲突 warning 或多列表排序。

\inst{django__django-14730} 是自引用对称 ManyToManyField 上 \code{related_name} 无效但缺少警告的问题。委员会补测生成的 patch 修改 \code{django/db/models/fields/related.py} 和 \code{docs/ref/checks.txt}，新增 \code{fields.W345} 警告，并记录用户传入的 \code{related_name}，防止后续 \code{contribute_to_class()} 自动改写字段后丢失原始意图。这个样本里的强证据包括“自引用对称关系中 \code{related_name} 没有效果”和“Django check framework 应以 warning 形式提示无效选项”。弱证据则是一些人工构造的边界 probe，例如未绑定 field 的中间状态。委员会阶段明确区分运行时检查、文档条目、已有 warning 的回归保护和自引用边界，最终官方 F2P 与 P2P 都为 0。

这两个通过样本有一个共同点。它们要解决的都不是“缺一行代码”式的问题，而是要把行为要求、相邻测试、文档和 warning 体系一起对齐。委员会结构在这类样本上更容易发挥作用，因为它把“该检查什么”提前到了修复者写代码之前。

\section{失败与未评测样本}

已评测但未通过的 8 个样本有一个共同特征。它们的 patch 都能 apply，P2P 回归都是 0，失败全部来自 F2P 残留。这说明委员会 V1 没有把问题卡在 patch 格式或回归破坏上，真正的缺口是语义修复没有完全覆盖隐藏测试行为。

\inst{django__django-13321} 的残留失败最重，F2P 失败达到 18 个，说明补丁方向与隐藏测试契约距离较大。\inst{django__django-14155} 残留 3 个失败，\inst{django__django-11564}、\inst{django__django-11630} 和 \inst{django__django-11905} 各残留 2 个失败。这类样本更像“已经修到一部分，但少了关键边界”。\inst{astropy__astropy-14365}、\inst{astropy__astropy-7746} 和 \inst{django__django-13768} 各残留 1 个失败，它们最适合留给 V2 的补丁攻击阶段继续研究，看分支攻击、回归扫描和可见行为审查能不能找出最后那个 corner case。

未评测的 3 个样本按保守口径计为未解决。\inst{django__django-11283} 和 \inst{django__django-15695} 没有形成干净非空 patch。\inst{django__django-15252} 的干净委员会补测没有形成正式非空 patch。它先前存在过一个辅助候选并完成过官方评测，但那属于辅助材料，最终统计仍按干净委员会补测口径处理。

\section{对验证数据的启发}

委员会实验指向一个明确方向。剩余失败大多不是补丁无法应用，也很少是 P2P 回归，真正缺的是更强的验证数据。这里的“更强”必须守住盲验证边界，既不去碰隐藏测试，也不把失败题的答案预先写给模型。可行的做法是让工作流在看到每一道题之后，自己从公开材料里抽取契约、相邻公开测试、可见行为字段和风险边界，再把这些材料交给裁判筛选。

V1 已经证明结构化验证能带来增量，13 个剩余样本里通过了 2 个。V2 则把下一步的研究问题写得更清楚。在没有官方隐藏测试信息的前提下，增加可见行为清单、契约仲裁、分支攻击、回归扫描和最终怀疑者，能不能把“只差一个边界”的样本转化为通过样本。如果后续继续实验，最值得观察的指标包括强证据临时检查的数量、弱证据诊断的数量、被最终怀疑者拦下的证据不足字段，以及每个未通过样本的首个 F2P 残留能否被 V2 提前预警。

\section{可复查文件}

\begin{center}
{\small
\setlength{\tabcolsep}{3pt}
\begin{longtable}{L{7.0cm} L{7.2cm}}
\toprule
文件或目录 & 用途 \\
\midrule
\endfirsthead
\toprule
文件或目录 & 用途 \\
\midrule
\endhead
\code{workflows/validation_council_solve.py} & 委员会 V1 的实际实现，包含定位、证据、修复前验证、修复、补丁风险审计和最终验证。 \\
\code{workflows/swe_committee_v2.py} & 委员会 V2 的图同构实现，包含可见行为清单、契约仲裁庭、补丁攻击阶段和最终怀疑者。 \\
\code{swebench/gen_prediction_workflow.py} & SWE-bench 工作流驱动，负责盲验证保护、预算、容器内生成和 guarded staged diff 抽取。 \\
\code{eval_work/swe_committee_mode_remaining13_glm52_20260629} & 委员会 V1 剩余 13 题主批次生成、metrics、predictions 和分批官方评测。 \\
\code{eval_work/swe_committee_targeted_12_34_8m_glm52_20260630} & 针对性补跑和官方评测材料，包含 \inst{django__django-11019} 的通过结果。 \\
\code{committee_mode_results.tex} & 较早的委员会结果小结，保留逐题结果和截止口径，位于本报告同一目录。 \\
\code{committee_report.tex} & 当前独立委员会报告源码，位于本报告同一目录。 \\
\bottomrule
\caption{委员会报告的主要可复查材料}
\end{longtable}
}
\end{center}

\section{结论}

委员会模式把 SWE-bench 解题从“让模型直接修”推进到“先建立证据，再写补丁，再攻击补丁，最后放行”的可审计结构。V1 已经在剩余 13 个难样本上新增解决 2 个，使前 100 题保守累计结果达到 89/100。失败样本的共同形态是 patch 可应用、P2P 为 0、F2P 仍有残留，说明后续提升重点应放在更强的盲验证和边界探索上。

V2 已经把这一方向拆成更细的组织结构，包含三路证据、契约仲裁、前置验证、三路补丁攻击、后置裁判、最终怀疑者和最终验证者。下一轮实验可以把 V2 单独作为新条件，与 V1 和 Team 模式分开统计，重点看它能不能减少“只差一两个 corner case”的语义失败。

\end{document}
